Apesar de talvez não ter notado, desde o seu primeiro código você vem lidando 
com funções. Funções nada mais são do que blocos de códigos capazes de 
executar uma sequência específica de instruções e são utilizadas para realizar 
uma melhor modularização do código, ou seja, separar o código ``em pedaços'' 
com diferentes funções bem definidas.
\\\\
As funções em C geralmente seguem o seguinte padrão de declaração:

\begin{verbatim}
    tipo_do_retorno nome_da_função(argumentos)
    {
    // Código da função
    }
\end{verbatim}

O tipo do retorno identifica o tipo de valor que a função deve retornar, 
caso necessário. Este valor pode ser de qualquer tipo já apresentado até o 
momento, como inteiros, caracteres, ponto flutuante, entre outros.
\\\\
Vamos considerar um exemplo simples para melhor entender a declaração de uma função:

\begin{verbatim}
    #include <stdio.h>

    void printHello(void) // Define uma nova função
    {
        printf("Hello!\n");
    }

    int main(void) // Define a função principal
    {
        printHello();
        printHello();
        printHello();
        return 0;
    }
\end{verbatim}

Este código é um exemplo básico de utilização de função. Ele define uma 
função chamada \textbf{printHello}, responsável por imprimir o texto 
"Hello!" três vezes consecutivas, sem necessidade de receber ou retornar 
qualquer valor.

\begin{ps}
    Um detalhe interessante sobre funções em C é que, devido ao fato de 
    a linguagem não depender de indentação, é possível declarar funções 
    em apenas uma linha, independentemente do seu tamanho.
    \\\\
    Isso não significa necessariamente que você deve fazer isso em todos 
    os casos, mas pode ser útil em algumas situações específicas. Se não 
    ficou claro como isso é possível, aqui está um exemplo utilizando a 
    mesma função \textbf{printHello}:

    \begin{verbatim}
    void printHello(void) { printf("Hello!\n"); }
    \end{verbatim}

    Se você substituir o código anterior por essa versão de apenas uma 
    linha, o código continuará sendo compilado normalmente.
\end{ps}

\subsubsection{Recebendo e retornando Valores}

As funções podem retornar valores após a execução de suas operações. Isso 
permite que elas comuniquem resultados de volta para o código que as 
chamou. Vamos elaborar outro exemplo para demonstrar isso:

\begin{verbatim}
    #include <stdio.h>

    // Define uma nova função que retorna 
    // o valor da soma de dois valores
    int soma(int a, int b) 
    {
        return a + b;
    }

    int main(void)
    {
        int resultado = soma(3, 4);
        printf("A soma é: %d\n", resultado);
        return 0;
    }
\end{verbatim}

Neste exemplo, a função \texttt{soma} recebe dois argumentos inteiros e 
retorna a soma deles. No \texttt{main}, chamamos essa função com os 
argumentos \(3\) e \(4\), atribuindo o resultado à sua variável, 
que é então impressa na tela. É importante observar que os argumentos devem 
ser inseridos na mesma ordem em que foram declarados na função.
\\\\
Os conceitos acima explicados são fundamentais para entender o funcionamento 
das funções em C e como elas podem ser utilizadas para tornar o código mais 
modular e eficiente. Ademais, o recomendado é que se utilize de bastante prática 
testes para ter certeza que você entendeu tudo que foi explicado.

\subsubsection{Incluindo funções de outros arquivos}

\subsubsection{Exercícios}

\textbf{É de extrema importância que você realize os exercícios aqui elaborados para fixar o 
conteúdo.}

\begin{enumerate}
    \item{Modifique os valores passados para a função soma para gerar diferentes resultados.}
    \item{Tente alterar os valores passados para a função soma para caracteres ou símbolos (`c', `\$' e etc.) e analise o comportamento do código.}
    \item{Modifique a função de maneira que seja possível somar 3 valores ao invés de 2.}
    \item{Elabore mais 3 funções para que seja possível realizar operações de subtração, multiplicação, divisão e tente utilizá-las ao invés da função soma.}
\end{enumerate}
